\section{Finite coefficient construction and rigorous remainder}\label{spm:sec:allorders}
Set $h=n^{-1/2}$ and make the central substitutions
\begin{equation}\label{spm:eq:scaling}
 t=h^{-2}(1+hy),\quad \theta=hw,\quad
 \rho(h)=\sqrt{1+v^2h^2/4}-vh/2,
\end{equation}
so that
\[
 x=h^{-1}\rho(h)e^{\ii hw},\qquad
 k=h^{-2}(1+hy)/L.
\]
Let $\mathcal E$ be the product Gaussian functional with independent variables $Y,W$ satisfying
\[
 Y\sim N(0,1),\qquad
 \PP(W\in\dd w)=\pi^{-1/2}e^{-w^2}\,\dd w.
\]
It is enough to define it on polynomials:
\[
 \mathcal E(Y^{2a})=(2a-1)!!,\qquad
 \mathcal E(W^{2b})=\frac{(2b-1)!!}{2^b},
\]
with zeroth moments $1$, odd moments $0$, and products factorized.

Define formal series with removable singularities at $h=0$ by
\begin{align}
 G(h,y)&=h^{-2}\bigl(\log(1+hy)-hy\bigr)+y^2/2,\label{spm:eq:G}\\
 H(h,w)&=v h^{-1}\rho(h)(e^{\ii hw}-1)
 +\frac{\rho(h)^2}{2h^2}(e^{2\ii hw}-1)-\frac{\ii w}{h}+w^2,\label{spm:eq:H}\\
 \mathcal B(h,y,w)&=B\bigl(h^{-1}\rho(h)e^{\ii hw},
                       h^{-2}(1+hy)/L,v\bigr).\label{spm:eq:scaledB}
\end{align}
Here $G(0,y)=H(0,w)=0$ and $\mathcal B(0,y,w)=S(L,v)$.

\begin{proposition}[Finite Gaussian algorithm]\label{spm:prop:algorithm}
The coefficients in Theorem~\ref{spm:thm:main} are defined by
\begin{equation}\label{spm:eq:algorithm}
 \sum_{j\geq0}C_j(L,v)h^j
 =e^{-S(L,v)}
 \frac{\mathcal E\exp(G+H+\mathcal B)}
      {(\mathcal E\exp G)(\mathcal E\exp H)}.
\end{equation}
This is an identity of formal power series: at any fixed order only finitely many Gaussian moments are taken. Both denominator factors are required.
\end{proposition}
\begin{proof}
On the central window, the gamma exponent is $-y^2/2+G$ and the angular exponent is $-w^2+H$. Dividing by the exact gamma normalization and the exact involution coefficient integral gives respectively the two factors in the denominator. In particular the gamma density is centered at $n$, not its mean $n+1$; the series $G$ and its normalizer retain that shift.

The least powers of $h$ in the three summand families of \eqref{spm:eq:B} are respectively $m-2$, $2m-4$, and $2m-2$. To calculate through $h^J$ one therefore needs only
\[
 m\leq J+2,\qquad m\leq\lfloor J/2\rfloor+2,
 \qquad m\leq\lfloor J/2\rfloor+1
\]
in those families. Each surviving coefficient is a polynomial in $y,w,L,v$ with rational coefficients apart from powers of $\ii$. Gaussian evaluation kills odd powers of $w$; the surviving powers of $\ii$ are real. Denominator division, with both constant terms equal to $1$, proves the rational-polynomial assertion.
\end{proof}

For transparency, the first calculation can be displayed without a symbolic system. Put $a=(v^2-1)/2$. The first nonconstant terms are
\begin{align*}
 G_1&=y^3/3,\qquad H_1=vw^2/2-2\ii w^3/3,\\
 \mathcal B_1&=aL(-v+2\ii w-y)
       +\frac{L^2}{4}(-2v+4\ii w-2y)+\frac{v^3L^2}{3}.
\end{align*}
The normalizers cancel the expectations of $G_1,H_1$, leaving
$C_1=\mathcal E\mathcal B_1$, exactly \eqref{spm:eq:C1}. At second order,
\begin{equation}\label{spm:eq:secondalgorithm}
 C_2=\mathcal E\left(\mathcal B_2+(G_1+H_1)\mathcal B_1
                         +\frac{\mathcal B_1^2}{2}\right)
       -\mathcal E(G_1+H_1)C_1.
\end{equation}
Substitution in this finite polynomial identity gives \eqref{spm:eq:D2}--\eqref{spm:eq:C2}. The reproduction package checks these identities exactly and separately checks the collision first correction.

\subsection{Why the formal algorithm is asymptotic}
A formal Gaussian calculation alone would not prove the theorem. Choose a fixed $0<\alpha<1/3$ and restrict $|y|,|w|\leq h^{-\alpha}$. The previous section bounds the excluded tails by $O(e^{-c h^{-2\alpha}})$. On this window, $1+hy$ is separated from zero, $\rho$ is analytic, and $|x/k|=O(h)$. The logarithmic series \eqref{spm:eq:B} and any fixed number of its $h$ derivatives converge uniformly.

Taylor's theorem through order $J$ for the local exponential has an integrable remainder bounded by
\begin{equation}\label{spm:eq:taylorbound}
 C_Jh^{J+1}(1+|y|+|w|)^{D_J}e^{\eta(y^2+w^2)},
 \qquad \eta<1/4.
\end{equation}
Here is a direct justification valid also at the intermediate points in Taylor's integral remainder. For every $0\leq s\leq h$, the identity
$\rho(s)^2+vs\rho(s)=1$ removes the apparent singularity in $H$. Uniformly on the window,
\[
 |G(s,y)|\leq Cs|y|^3,\qquad
 |H(s,w)|\leq Cs(|w|^2+|w|^3),\qquad
 |\mathcal B(s,y,w)|\leq C.
\]
Because $h(|y|+|w|)=o(1)$, the residual real parts are at most
$\eta(y^2+w^2)+C$ after taking $h$ sufficiently small. Every fixed derivative is bounded by a polynomial in $|y|+|w|$: the only denominator factors are powers of $1+sy$ and analytic factors separated from zero, and the least-degree structure of \eqref{spm:eq:B} bounds the number of relevant terms. Differentiating the exponential a fixed number of times multiplies it only by such polynomials. This proves \eqref{spm:eq:taylorbound}.

Multiplication by the Gaussian density $e^{-y^2/2-w^2}$ makes \eqref{spm:eq:taylorbound} integrable over the full plane. Its integral is $O_J(h^{J+1})$. Extending each finite polynomial coefficient integral from the window to the full plane loses only an exponentially small tail. The same argument applies to both denominator integrals, whose constant terms are nonzero. This proves \eqref{spm:eq:main} after the exponentially small pole error \eqref{spm:eq:poleerror} is restored.

The notation $\mathcal E e^{G+H+\mathcal B}$ in \eqref{spm:eq:algorithm} is formal. The proof extends only finite polynomial truncations to the full real plane; it never integrates $\log(1+hy)$ through $y=-1/h$.

Finally the normalized quantity in \eqref{spm:eq:main} is holomorphic and nonzero on a slightly smaller neighborhood for large $n$, by Lemma~\ref{spm:lem:involution} and its uniform limit $1$. Cauchy's formula, applied on a larger compact polydisc, bounds all fixed derivatives of the remainder. Its analytic logarithm consequently has a differentiable expansion to every fixed order. This completes the proof of Theorem~\ref{spm:thm:main}.

\begin{remark}[{[write]} The third coefficient at $u=v=1$, and the data, 6~October 2026]\label{spm:rem:third}
The write evaluated the finite prescription~\eqref{spm:eq:algorithm} at
$u=v=1$ through $h^3$ in SymPy~1.14.0, with $L$ symbolic (an implementation
of its own, not shipped; it expands $G$, $H$ and $\mathcal B$ from
\eqref{spm:eq:G}--\eqref{spm:eq:scaledB}, keeping in~\eqref{spm:eq:B}
exactly the terms listed in the proof of
Proposition~\ref{spm:prop:algorithm}, takes the Gaussian moments of
$\mathcal E$, and divides by both normalizers). It returns $S(L,1)=L^2/4$,
$C_1(L,1)=-L^2/6$ and $C_2(L,1)=-L^2/4+L^4/18$, as in
\eqref{spm:eq:C1}--\eqref{spm:eq:C2}, and
\[
 C_3(L,1)=-\frac{L^2(50L^4+324L^2-1755)}{6480},\qquad
 D_3(L,1)=C_3-C_1C_2+\tfrac13C_1^3=-\frac{L^2(22L^2-65)}{240}.
\]
Theorem~\ref{spm:thm:main} with $J=3$ therefore extends~\eqref{spm:eq:unmarked}
by the term $C_3n^{-3/2}$, with remainder $O(n^{-2})$; at $L=\log2$,
$C_1=-0.08007550\ldots$, $C_2=-0.10728908\ldots$ and
$C_3=0.11772518\ldots$. This value rests on the write's single
implementation (Question~\ref{spm:q:third}).

\emph{The data.} With $\varrho_n=A_n\big/\bigl(e^{L^2/4}I_n(1)/(2L^{n+1})\bigr)$
and $h=n^{-1/2}$, the exact counts of Remark~\ref{spm:rem:oeis} give
(50-digit \texttt{mpmath}; the delivered noncertified receipt
\file{data/diagnostic_receipt.json} records the quantities of the second to fifth columns at
$n=32$, $80$, $160$, $320$, $640$, and agrees with the row $n=32$):
\begin{center}\small
\begin{tabular}{@{}r*{5}{c}@{}}
\toprule
$n$ & $\varrho_n$ & $(\varrho_n-1)/h$ & $(\varrho_n-1-C_1h)/h^2$ & $(\varrho_n-1-C_1h-C_2h^2)/h^3$ & $(\varrho_n-\sum_{j\le3}C_jh^j)/h^4$\\
\midrule
16 & 0.97438 & $-0.1025$ & $-0.0896$ & 0.0709 & $-0.1872$\\
32 & 0.98295 & $-0.0965$ & $-0.0927$ & 0.0826 & $-0.1986$\\
50 & 0.98678 & $-0.0935$ & $-0.0947$ & 0.0889 & $-0.2041$\\
100 & 0.99102 & $-0.0898$ & $-0.0976$ & 0.0966 & $-0.2112$\\
200 & 0.99384 & $-0.0872$ & $-0.1000$ & 0.1024 & $-0.2168$\\
300 & 0.99504 & $-0.0859$ & $-0.1012$ & 0.1051 & $-0.2194$\\
400 & 0.99574 & $-0.0852$ & $-0.1020$ & 0.1067 & $-0.2210$\\
500 & 0.99621 & $-0.0847$ & $-0.1025$ & 0.1078 & $-0.2222$\\
\bottomrule
\end{tabular}
\end{center}
(rounded to the digits shown). The third to fifth columns move toward
$C_1$, $C_2$ and $C_3$ as $n$ grows, and the last stays between $-0.18$
and $-0.23$ for $16\le n\le500$. These are observations on a finite range,
consistent with the expansion; they certify no coefficient and no rate.
\end{remark}
\begin{writenote}[Independent check of the write, 6~October 2026]
An adversarial check made by the intake after the write (commit
\file{c569b0b98}) re-derived the statements marked [write] with its own
code, after fetching again A138178, its b-file and A135588 (the revisions
quoted) and Cameron--Prellberg--Stark's v2 (4~April 2006).
\emph{Counts.} By a route different from the write's, the zero-line
inclusion--exclusion~\eqref{spm:eq:deletion} with $f_n(j)$ from the
first-order recurrence
$(n+1)a_{n+1}=ja_n+(n-1+j+j(j-1))a_{n-1}$ of
$(1-z)^{-j}(1-z^2)^{-j(j-1)/2}$, it computed $A_0,\ldots,A_{2000}$ exactly:
all 501 b-file terms and the 25 data terms agree; binomial cell factors
give all 26 terms of A135588. \emph{Remark~\ref{spm:rem:third}.} Every entry
of the table was reproduced to the digits shown, and the receipt's row
$n=32$ agrees. Richardson extrapolation in $h$ of
$(\varrho_n-1-C_1h-C_2h^2)/h^3$ from the exact counts gives
$0.1177251847$ (nodes $n=400,600,\ldots,2000$; the same ten digits from
three other node sets of seven to thirteen points), against $C_3(\log2,1)=0.1177251846518\ldots$ from the
write's formula; the same extrapolation of marked counts $A_n(u,1)$ to
$n=1500$ gives $0.154285415$ at $u=\tfrac45$ and $0.087284579$ at
$u=\tfrac54$, against $0.1542854152\ldots$ and $0.0872845791\ldots$ from the
formula at $L=\log\tfrac94$ and $L=\log\tfrac95$. This is numerical
confirmation at three values of $L$, independent of the Gaussian
prescription; it is not a derivation (Question~\ref{spm:q:third}).
$D_3=C_3-C_1C_2+\tfrac13C_1^3$ and $C_1(L,1)$, $C_2(L,1)$ were rechecked in
SymPy. \emph{Remark~\ref{spm:rem:oeis}.} The entries' revisions, names,
generating functions (the first is~\eqref{spm:eq:deletion} at $u=v=1$,
Jovovic's is~\eqref{spm:eq:transform}), comments and the absence of an
asymptotic are as quoted; Proposition~4.2 (p.~14, with $C_s\approx0.44341$)
and the Poisson$((\log2)^2/2)$ limit with~(8) (p.~9) are in the preprint as
described. \emph{Section~\ref{spm:sec:model}, note.} \file{q2:thm:fubini},
Part~IV of \file{a261781-matrix-compositions} (Sections~34--44, the polydisc
$|\omega-1|\le2.1$, $|u-1|\le0.02$, $P\sim\mathrm{Pois}(Lr)$ with
$r\log2=1.1046161\ldots$) and the Lean names are as stated.
\emph{Section~\ref{spm:sec:saddles}, note.} The identity in law was
re-derived ($\sum_{k\ge d}q^k\binom kd=u^d/(1-q)$, and $K$ has the law
proportional to $e^{-Lk}f_n(k,v)=q^kf_n(k,v)$ there).
\emph{Remark~\ref{spm:rem:transseries}.} (1)~holds; (2)~the factorial core
with $\kappa=\tfrac12$, $d=-\tfrac12-\log L$ gives $v_{\rm vol}=w$,
$X_{\rm vol}=N$, slope $Q/2$, and the core equation alone has the
\file{plt:thm:lw-template} data stated, the extra terms being of relative
order $e^{-Z/2}/Z$; (3)~the identity for $F_0$ is exact, the triangular
recurrence gives $e_1,e_2,e_3=\alpha,\beta,\gamma$, and
$e_1,\dots,e_5$ are polynomials in $Q^{-1}$ without constant term, of
degrees $1,3,5,7,9$, with $d$ first in $e_4$ (SymPy); (4)~is accurate.
\emph{Provenance.} The archive facts (604{,}411 bytes, the SHA-256 quoted,
30 files, 29 manifest entries, 952 delivered lines, 22 pages) and the counts
of 110 labels and 93 references were confirmed. No defect was found.
\end{writenote}
